\section{Experiments and Results}
\label{sec:results}

\subsection{Experimental Setup}

After cold-start filtering (at least 5 interactions per user and 3 users per
song) the modelling dataset contains \TODO{$N_\text{users}$} users
($\approx$284{,}864) and \TODO{$N_\text{songs}$} tracks ($\approx$7{,}050) with
their full jSymbolic feature vectors and knowledge-graph context. Interactions
are partitioned with the user-level stratified split of
\cref{sec:eval-protocol} (\TODO{report the train/val/test interaction counts}).
All models are trained on identical training interactions and evaluated on the
held-out test interactions; reproducibility is fixed with seed~42. Hyper\-
parameters are those listed in \cref{sec:methodology} and were selected on the
validation split only.

\subsection{Quantitative Comparison}

\cref{tab:main-results} reports every model at $K=10$ on the test set. The
popularity prior is, as expected, a strong accuracy baseline but collapses on
diversity (it reaches only 0.6\% catalogue coverage and the highest popularity
bias). KNN-CF trades a modest amount of ranking accuracy for a dramatic gain in
coverage and a large drop in popularity bias, and consequently wins on the
composite \emph{Overall} score. \TODO{Update the XGBoost-Hybrid and HGT rows
with final numbers once the autoencoder embeddings have been regenerated;
state which model achieves the best Overall score and by how much.}

\begin{table}[t]
\centering
\caption{Top-10 test-set performance. Higher is better except
PopBias@10 ($\downarrow$). Baseline rows are measured; graph/hybrid rows are
pending the regenerated embeddings.}
\label{tab:main-results}
\footnotesize
\setlength{\tabcolsep}{3.5pt}
\begin{tabular}{l S[table-format=1.4] S[table-format=1.4] S[table-format=1.4] S[table-format=1.4] S[table-format=1.4]}
\toprule
\textbf{Model} & {Recall} & {NDCG} & {HitRate} & {Cov.} & {Overall} \\
\midrule
MostPopular           & 0.1285 & 0.1143 & 0.2002 & 0.0064 & 0.1477 \\
KNN-CF ($k$=\TODO{5}) & 0.0923 & 0.0971 & 0.1540 & 0.9176 & 0.4185 \\
XGBoost-Hybrid        & \TODO{} & \TODO{} & \TODO{} & \TODO{} & \TODO{} \\
HGT (ours)            & \TODO{} & \TODO{} & \TODO{} & \TODO{} & \TODO{} \\
\bottomrule
\end{tabular}
\end{table}

\paragraph{Behaviour across list length.}
\cref{fig:multik} sweeps the recommendation-list length $K\in\{5,\dots,100\}$
for Recall@$K$, NDCG@$K$ and the composite Overall score, exposing the
``sweet-spot'' list length for each model. \TODO{Insert the multi-$K$
comparison figure produced by the baseline-comparison cell and comment on where
each model peaks.}

\begin{figure}[t]
\centering
\fbox{\parbox[c][3.2cm][c]{0.92\linewidth}{\centering
  \TODO{Insert \texttt{Recall@K / NDCG@K / Overall} vs.\ $K$ line plots
  (baseline-comparison cell).}}}
\caption{Metric vs.\ recommendation-list length $K$ for all models.}
\label{fig:multik}
\end{figure}

\paragraph{Statistical significance.}
Pairwise paired Wilcoxon signed-rank tests over all
$\approx$284{,}864 test users find the accuracy differences between models to be
statistically significant ($p<0.05$), but with small Cohen's~$d_z$ effect sizes
($\approx$0.03--0.11), confirming that the practically meaningful separation
lies in coverage and popularity bias rather than raw ranking accuracy. An
omnibus Friedman test across all models is reported in \cref{tab:friedman}.
\TODO{Populate \cref{tab:friedman} (and the pairwise table) from the comparison
cell once all four models are final.}

\begin{table}[t]
\centering
\caption{Friedman omnibus test and best-ranked model per metric.}
\label{tab:friedman}
\footnotesize
\begin{tabular}{l c c l}
\toprule
\textbf{Metric} & \textbf{$\chi^2$} & \textbf{$p$} & \textbf{Best model} \\
\midrule
Recall@$K$  & \TODO{} & \TODO{} & \TODO{} \\
NDCG@$K$    & \TODO{} & \TODO{} & \TODO{} \\
HitRate@$K$ & \TODO{} & \TODO{} & \TODO{} \\
MRR         & \TODO{} & \TODO{} & \TODO{} \\
\bottomrule
\end{tabular}
\end{table}

\subsection{Learning Curves}

\cref{fig:curves} collects the training-time diagnostics for the trained
models: autoencoder reconstruction MSE (train/val), the RotatE/ComplEx KGE
training loss, the HGT debiased-listwise loss together with validation Recall
and NDCG, and the XGBoost per-round \texttt{ndcg@10}. \TODO{Drop in the four
PNGs exported to \texttt{data/final/} (\texttt{ae\_loss\_curve.png},
\texttt{kge\_loss\_curve.png}, \texttt{hgt\_training\_curves.png},
\texttt{xgb\_loss\_curve.png}); comment on convergence and any over-fitting.}

\begin{figure}[t]
\centering
\fbox{\parbox[c][4cm][c]{0.92\linewidth}{\centering
  \TODO{Insert learning curves: AE (MSE), KGE (loss), HGT (loss + val
  Recall/NDCG), XGBoost (ndcg@10 per round).}}}
\caption{Training and validation learning curves for the trained models.}
\label{fig:curves}
\end{figure}

\subsection{Qualitative Analysis}

We complement the metrics with three qualitative views.
\textbf{(i)~Population level:} per-model distributions of cosine-to-profile
similarity and Jensen--Shannon divergence between the genre/mode/tempo/decade
profiles of a user's training history and their recommendations, plus
dominant-attribute match rates (\cref{fig:qual}).
\textbf{(ii)~Per-user case studies:} successful users across different
training-history sizes (light/medium/heavy), failure cases, and ``mixed'' users
where models disagree, each shown with a cross-model metric table and
train-vs-recommendation attribute distributions.
\textbf{(iii)~Latent space:} a UMAP/GMM projection of the learned node
embeddings used to inspect cluster structure.
\TODO{Insert the population-summary mosaic, one or two representative per-user
case figures, and the latent-space projection; summarise where each model is
strong or weak (e.g.\ popularity vs.\ niche taste, cold vs.\ rich profiles).}

\begin{figure}[t]
\centering
\fbox{\parbox[c][3.6cm][c]{0.92\linewidth}{\centering
  \TODO{Insert population-qualitative mosaic and a per-user case study.}}}
\caption{Qualitative comparison of recommendation behaviour across models.}
\label{fig:qual}
\end{figure}

\subsection{Discussion}

The measured baselines already illustrate the central tension of the task:
popularity maximises short-list accuracy while almost entirely ignoring the
catalogue, whereas content/structure-aware models spread recommendations far
more widely at a small accuracy cost. \TODO{Once the HGT and hybrid numbers are
in, state explicitly whether adding the knowledge graph and musical features
improves the accuracy/diversity trade-off over the interaction-only baselines,
and connect back to the cold-start and popularity-bias motivation of
\cref{sec:intro}.}
